% !TeX root = ../main.tex
\section{Introduction}
\label{sec:introduccion}

The rapid proliferation of Artificial Intelligence (AI) and Machine Learning (ML) has transformed critical sectors, including healthcare, finance, and intelligent transportation systems. However, training high-capacity models requires massive volumes of data, which frequently contain sensitive or proprietary information, such as high-resolution geolocation traces, user behavior patterns, or biometric identifiers. Under traditional centralized paradigms, gathering raw data into a single repository introduces severe privacy vulnerabilities, regulatory compliance challenges (such as the GDPR), and substantial computational and communication bottlenecks.

Federated Learning (FL)~\cite{mcmahan2017federated} emerged as a decentralized paradigm where distributed edge nodes iteratively train local model instances on their private datasets, transmitting only updated weight vectors or gradients to a central coordinator. Although FL avoids direct exposure of raw data, it introduces a structural security challenge: the absence of inherent trust among participating entities. In an unsupervised, distributed environment, malicious actors or compromised edge nodes can launch data poisoning, model poisoning, or backdoor insertion attacks~\cite{bagdasaryan2020backdoor} by submitting manipulated model updates. For instance, in Connected and Autonomous Vehicles (CAVs)~\cite{chellapandi2024federated}, a subtle, coordinated perturbation in model parameters could lead the global classifier to misidentify critical road hazards or traffic signals (e.g., failing to detect a STOP sign), yielding catastrophic consequences.

To address these vulnerabilities, integrating Blockchain technology with Federated Learning (BCFL)~\cite{qu2022blockchain} provides an immutable, decentralized, and verifiable ledger~\cite{nakamoto2008bitcoin, zheng2018blockchain}. By embedding a distributed ledger into the federated training workflow, parameter updates, aggregation rounds, and participant identities are cryptographically anchored via smart contracts, establishing an end-to-end, auditable, and non-repudiable record of all training contributions.

Nevertheless, this convergence introduces a critical system-level trade-off: \textbf{security versus operational performance}. While blockchain provides immutability and Byzantine fault tolerance, it inherently introduces network latency, consensus delays, transaction confirmation overhead, and additional computational resource demands. In mission-critical edge computing environments~\cite{liu2019edge}, it is essential to answer a fundamental question: \textit{what temporal and computational overhead must be incurred to achieve a more auditable and tamper-resistant federated training?}

\subsection{Motivation and Prior Work}
This work builds upon and extends two research lines conducted at the University of Alicante:
\begin{enumerate}
    \item \textbf{Real-time vehicular event verification}~\cite{pujol2024blockchain}: demonstrated the feasibility of a private Hyperledger Besu network operating under Istanbul Byzantine Fault Tolerant (IBFT 2.0) consensus to log vehicular traffic alerts with sub-second latencies (0.32 to 0.82~s), outperforming human reaction times.
    \item \textbf{FL-Blockchain architectural framework}~\cite{cloudlabaia_flblockchain}: conceptually proposed using smart contracts to record cryptographic hashes of model weights at each federated round.
\end{enumerate}

The present study formally designs, implements, and empirically validates the convergence of these proposals on a physical distributed testbed, evaluating both cryptographic integrity guarantees and temporal system performance.

\subsection{Main Contributions}
The main contributions of this paper are summarized as follows:
\begin{itemize}
    \item \textbf{Hybrid architecture design}: We formulate an end-to-end framework coupling a permissioned Hyperledger Besu blockchain network with the Flower federated learning framework~\cite{beutel2020flower} and off-chain, content-addressable distributed storage via IPFS.
    \item \textbf{Traceability and integrity via smart contracts}: We design and deploy a Solidity smart contract that orchestrates bidirectional cross-verification of models through content identifiers (IPFS CIDs, encapsulating SHA-256 multihashes of serialized artifacts), immutably anchoring each participant's contribution and the aggregated global model per round.
    \item \textbf{Empirical adversarial validation}: We subject the platform to deliberate weight-corruption attack vectors (both in-transit tampering from an untrusted server and unauthorized/poisoning injections from clients), validating that the bidirectional verification protocol deterministically intercepts integrity discrepancies and rogue models with low response latencies (31 to 149~ms).
    \item \textbf{Ex-post audit protocol and mathematical FedAvg verification}: We develop and evaluate an independent, off-band audit protocol enabling any consortium participant or third-party auditor to retrospectively verify aggregator fidelity. We demonstrate the ability to mathematically reconstruct canonical FedAvg aggregation directly from immutable commitments stored on Besu and IPFS, achieving exact numerical correspondence ($\Delta = 0$, CID match) in $0.34$~s, while detecting unauthorized parameter deviations in $0.18$~s.
    \item \textbf{Quantitative infrastructure benchmarking}: We deploy a convolutional traffic event classifier to serve as a representative workload for extracting reproducible system-level metrics (temporal overhead, EVM gas consumption, and validator scaling), delimiting hyperparameter and model optimization as future work.
    \item \textbf{Operational feasibility analysis}: We analytically resolve the latency-versus-security dilemma, characterizing the operational regimes where blockchain logging overhead is thoroughly justified compared to unverified setups.
\end{itemize}

To facilitate scientific reproducibility and technology transfer, the complete framework source code, Solidity smart contracts, Hyperledger Besu Docker Compose deployment templates, and experimental evaluation scripts are openly available on GitHub: \url{https://github.com/dpz6-ua/Blockchain_ML}.

\subsection{Paper Organization}
The remainder of this paper is structured as follows: Section~\ref{sec:estado_arte} reviews related work and theoretical foundations of Blockchain and Federated Learning. Section~\ref{sec:arquitectura} details the proposed architecture and smart contract design. Section~\ref{sec:implementacion} describes the experimental testbed and dataset. Section~\ref{sec:evaluacion} presents and discusses empirical performance, scalability, and security validation results. Finally, Section~\ref{sec:conclusiones} concludes the paper and outlines future research directions.
